\section*{Hoofdstuk \ref{ch:g7:sound-production} --- Geluid: opwekking en voortplanting}

\begin{solution}{exo:g7:sound-production:1}
De stof waar een \omterm{def:g3:sound-around-us:sound}{geluid} doorheen reist --- \omterm{def:g2:air-around-us:air}{lucht}, water, hout,
staal, elke moleculaire menigte. Het \omterm{def:g7:sound-production:medium}{niet-medium}: vacuüm.
\end{solution}

\begin{solution}{exo:g7:sound-production:2}
De estafette reist: een snellend patroon van samenpersingen dat van rij
tot rij door de menigte wordt doorgegeven. Elk \omterm{def:g7:states-of-matter:molecule}{molecuul} wiebelt
alleen om zijn eigen plaats --- de boodschap beweegt, de boodschappers
blijven.
\end{solution}

\begin{solution}{exo:g7:sound-production:3}
Geen menigte, geen estafette: in de leeggepompte stolp duwen de
trillende wanden van de bel tegen niets, dus vormt zich geen boodschap.
Ruimteslagen zijn om dezelfde reden stil --- ontploffingen tussen de
sterren hebben geen menigte om hun gebrul te dragen.
\end{solution}

\begin{solution}{exo:g7:sound-production:4}
\omterm{def:g2:air-around-us:air}{Lucht} ongeveer \qty{340}{m/s}; water ongeveer \qty{1500}{m/s};
staal ongeveer \qty{5000}{m/s}. Hoe strakker de \omterm{def:g7:states-of-matter:molecule}{moleculen} elkaar
raken, hoe sneller de duw wordt doorgegeven: vliegend, rakend,
vastgezet.
\end{solution}

\begin{solution}{exo:g7:sound-production:5}
De vastgezette rijen van staal geven bijna van hand tot hand door:
ongeveer vijftien keer het tempo van de \omterm{def:g2:air-around-us:air}{lucht}, dus komt de
boodschap van de rail lang voor het gerommel door de \omterm{def:g2:air-around-us:air}{lucht} aan.
\end{solution}

\begin{solution}{exo:g7:sound-production:6}
De reis van het \omterm{def:g1:light-and-shadows:source}{licht} over welk veld dan ook telt naast die van het
\omterm{def:g3:sound-around-us:sound}{geluid} als ogenblikkelijk: de geziene klap markeert het echte begin
met veel meer nauwkeurigheid dan de hand aan de stopwatch kan
benutten.
\end{solution}

\begin{solution}{exo:g7:sound-production:7}
$340 \times 1 = \qty{340}{m}$ heen en terug: de muur staat op
\qty{170}{m}. Wie \qty{340}{m} antwoordt, vergeet dat die \omterm{def:g2:measuring-time:second}{seconde} het
heengaan \emph{en} het terugkomen betaalde.
\end{solution}

\begin{solution}{exo:g7:sound-production:8}
In water: $1500 \times 2 = \qty{3000}{m}$ heen en terug --- diepte
\qty{1500}{m}. De misstap: de \qty{340}{m/s} van de \omterm{def:g2:air-around-us:air}{lucht} lenen
voor een reis onder water.
\end{solution}

\begin{solution}{exo:g7:sound-production:9}
$d = 340 \times 2.5 = \qty{850}{m}$. De flits van het vuurwerk gaf een
echt startschot; de ongeziene straaljager biedt er geen --- geen moment
van uitzending om vanaf te klokken --- en erger nog: het gehoorde
gebrul verliet de jager \omterm{def:g2:measuring-time:second}{seconden} geleden: het oor wijst naar waar de
jager \emph{was}.
\end{solution}

\begin{solution}{exo:g7:sound-production:10}
Het oor onder water wordt bediend door de estafette van water op
\qty{1500}{m/s}, het oor op de kant door die van de \omterm{def:g2:air-around-us:air}{lucht} op
\qty{340}{m/s} --- en een groot deel van het \omterm{def:g3:sound-around-us:sound}{geluid} van de
luidspreker steekt de grens aan het oppervlak überhaupt niet over naar
de \omterm{def:g2:air-around-us:air}{lucht}. Sneller \omterm{def:g7:sound-production:medium}{medium}, privéboodschap: de zwemmer wint
ondanks de afstand.
\end{solution}

\begin{solution}{exo:g7:sound-production:11}
Eén slag, twee estafettes: door staal en door \omterm{def:g2:air-around-us:air}{lucht}. Staal:
$3000 \div 5000 = \qty{0.6}{s}$; \omterm{def:g2:air-around-us:air}{lucht}: $3000 \div 340 \approx
\qty{8.8}{s}$ --- de klap loopt de dreun ongeveer \qty{8.2}{s} voor.
\end{solution}

\begin{solution}{exo:g7:sound-production:12}
Eén klap stuurt hetzelfde ogenblik beide media in; op \qty{750}{m} legt
de recorder de aankomst door water (op de hydrofoon) en de aankomst
door \omterm{def:g2:air-around-us:air}{lucht} (op de microfoon) vast. Verwachte tijden: water $750
\div 1500 = \qty{0.5}{s}$, \omterm{def:g2:air-around-us:air}{lucht} $750 \div 340 \approx
\qty{2.2}{s}$ --- een gat van ongeveer \qty{1.7}{s}. Achterstevoren:
met bekende $d$ en bekende \omterm{def:g7:motion-average-speed:formula}{snelheid} in \omterm{def:g2:air-around-us:air}{lucht} geeft het
gemeten gat de aankomsttijd door water ($t_{\text{lucht}} -
\text{gat}$), en $v = 750 \div 0.5 = \qty{1500}{m/s}$ --- het tempo van
water gemeten zonder ooit de stille klap zelf te klokken.
\end{solution}

\begin{solution}{pb:g7:sound-production:1}
\textbf{1.} $340 \times 4 = \qty{1360}{m}$ heen en terug: breedte
\qty{680}{m}.
\textbf{2.} $340 \times 1.5 = 510$; de helft: \qty{255}{m}.
\textbf{3.} Eerste wand: $340 \times 1 \div 2 = \qty{170}{m}$; tweede:
$340 \times 2 \div 2 = \qty{340}{m}$; de kloof overspant langs die lijn
$170 + 340 = \qty{510}{m}$.
\textbf{4.} De formule is exact; de hand niet: menselijk starten en
stoppen spreidt over tienden van een \omterm{def:g2:measuring-time:second}{seconde} --- tientallen
\omterm{def:g2:measuring-length:units}{meters} bij \qty{340}{m/s}. Vandaar de oude regel: herhaal, kijk hoe
de aflezingen samendrommen, en rapporteer de drom.
\textbf{5.} $1500 \times 0.2 = \qty{300}{m}$ heen en terug: diepte
\qty{150}{m}.
\textbf{6.} $1500 \times 0.4 \div 2 = \qty{300}{m}$.
\textbf{7.} Aan de grens lucht--water wordt het grootste deel van een
schreeuw teruggekaatst --- er komt weinig de rivier in, en de zwakke
bodemecho moet die grens opnieuw oversteken om een oor in de
\omterm{def:g2:air-around-us:air}{lucht} te bereiken: er overleeft vrijwel niets. En zelfs met die
\omterm{ex:g4:how-sound-travels:echo}{echo} verslaan reizen heen en terug van tienden van een \omterm{def:g2:measuring-time:second}{seconde} elke
tijdmeting op het gehoor.
\textbf{8.} De echotijd wordt abrupt langer wanneer de boot over de
rand \omterm{def:g1:floating-and-sinking:words}{drijft} --- diepere bodem, langere reis heen en terug: de sonar
tekent de klif onder water.
\textbf{9.} Staal: $2000 \div 5000 = \qty{0.4}{s}$; \omterm{def:g2:air-around-us:air}{lucht}: $2000
\div 340 \approx \qty{5.9}{s}$.
\textbf{10.} Beide estafettes starten samen, dus groeit het gat
evenredig met de afstand: elke kilometer geeft de trage estafette een
vaste extra vertraging op de snelle. Meet het gat, deel door het
verschil per kilometer, en de afstand rolt eruit --- één hamerslag,
zonder stopwatch aan het andere uiteinde.
\textbf{11.} Het railgeluid: mist is een menigte druppeltjes die
\emph{\omterm{def:g1:light-and-shadows:source}{licht}} hopeloos verstrooit maar een estafette die opgesloten
zit in vast staal nauwelijks hindert (ook luchtgeluid
overleeft mist goed --- maar de boodschap van de rail is de schoonste
en snelste van de drie).
\textbf{12.} Bijvoorbeeld: ``Elke \omterm{ex:g4:how-sound-travels:echo}{echo} betaalt de reis heen en
terug: halveer voor je hem gelooft. Elk \omterm{def:g7:sound-production:medium}{medium} bepaalt zijn eigen
tempo: kies de \omterm{def:g7:motion-average-speed:formula}{snelheid} bij de menigte, nooit uit gewoonte. En elke
tijdmeting in deze kloof begon met het oog --- de reis van het
\omterm{def:g1:light-and-shadows:source}{licht} telt hier als ogenblikkelijk, en daarom mag het zien van de
vallende hamer als startschot dienen.''
\end{solution}
